%=====================================================================
\chapter{CPU and Memory Virtualization}\label{ch:cpumem}
%=====================================================================
\locator{2}{Part II --- Hypervisors and Compute Virtualization}

\begin{outcomes}
By the end of this chapter you will be able to:
\begin{tight}
  \item \textbf{Explain} how a hypervisor schedules virtual processors onto
        physical ones, and \textbf{interpret} ready time as the measure of
        whether it is succeeding.
  \item \textbf{Explain} why giving a guest more virtual processors can make it
        slower, using co-scheduling.
  \item \textbf{Compute} CPU and memory overcommit ratios, and \textbf{decide}
        from measured figures whether a host can take another guest.
  \item \textbf{Order} the four memory reclamation techniques by cost and
        \textbf{state} what each one does to guest performance.
  \item \textbf{Describe} NUMA and \textbf{explain} why a guest sized larger
        than a socket may lose a third of its memory bandwidth.
\end{tight}
\end{outcomes}

\begin{prereq}
\chref{hardware} (extended page tables, which are what make guest memory
management cheap) and \chref{platforms} (the four schedulers). From Operating
Systems you need process scheduling and the idea of a working set.
\end{prereq}

\begin{keyterms}
vCPU $\bullet$ pCPU $\bullet$ time slice $\bullet$ ready time $\bullet$ \%RDY
$\bullet$ co-scheduling $\bullet$ co-stop $\bullet$ relaxed co-scheduling
$\bullet$ overcommit ratio $\bullet$ working set $\bullet$ active memory
$\bullet$ transparent page sharing $\bullet$ ballooning $\bullet$ balloon
driver $\bullet$ memory compression $\bullet$ hypervisor swapping $\bullet$
memory tax $\bullet$ NUMA $\bullet$ NUMA node $\bullet$ vNUMA $\bullet$ node
interleaving
\end{keyterms}

\section{Scheduling a Processor That Is Not There}

A guest configured with four virtual processors believes it has four
processors. It runs four run queues, balances threads across them, and spins
on locks expecting the other holder to be making progress. None of that is
true: a \term{vCPU} is a thread the hypervisor schedules onto a \term{pCPU},
and between one guest instruction and the next it may have been off the
processor for milliseconds.

For CPU-bound work this is nearly free --- the guest's instructions run on the
real processor at full speed, as \chref{architectures} measured. The costs
appear when guests contend, and they are measured by one number.

\begin{definitionbox}[Ready Time]
\term{Ready time} is the proportion of time a vCPU was \emph{runnable but not
scheduled} --- the guest had work, and no physical core was free for it.
VMware reports it as \term{\%RDY} per vCPU; KVM exposes the same quantity as
steal time, which the guest itself can see in \texttt{top} as \texttt{\%st}.

\smallskip
It is the single most useful number in capacity planning, because it
distinguishes the two cases that look identical from inside the guest: a guest
that is slow because its own work is slow, and a guest that is slow because it
is not being given a processor.

\smallskip
Rules of thumb, widely used and worth knowing: under 5\% is healthy, 5--10\%
deserves investigation, and sustained figures above 10\% mean the host is
oversubscribed. \chref{performance} treats the measurement properly.
\end{definitionbox}

\section{Why More vCPUs Can Be Slower}

The first instinct on finding a slow guest is to give it more processors. It
is often exactly wrong, and the reason is \term{co-scheduling}.

A guest kernel assumes its processors run simultaneously. A thread holding a
spinlock on vCPU~0 will be released shortly, so a thread on vCPU~1 spins rather
than sleeping. If vCPU~0 is \emph{descheduled} while holding that lock, vCPU~1
spins for a whole time slice --- burning a physical core to wait for a virtual
one that is not running.

Early hypervisors avoided this by strict co-scheduling: never run a guest's
vCPUs unless \emph{all} of them can run together. That removes the problem and
creates a worse one --- an eight-vCPU guest waits for eight free cores, and on a
busy host it waits a long time. Modern schedulers use \term{relaxed
co-scheduling}: a vCPU is allowed to fall behind its siblings by a small margin,
and is forced to catch up only when the skew grows too large. The waiting that
results is reported as \term{co-stop}.

\dsfig{%
\begin{tikzpicture}[font=\scriptsize,
  lab/.style={font=\scriptsize\bfseries, text=neutral!40!black, anchor=east},
  run/.style={draw=greenhl!70, fill=greenhl!25, line width=0.5pt},
  wait/.style={draw=accent!70, fill=accent!18, line width=0.5pt},
  note/.style={font=\scriptsize\itshape, align=left, anchor=west}
]
  % ---- 2-vCPU guest: fits easily ----
  \node[font=\scriptsize\bfseries, text=primary, anchor=west] at (0,3.5)
       {A 2-vCPU guest on a busy host};
  \foreach \i/\y in {0/3.0, 1/2.55}{
    \node[lab] at (-0.15,\y) {vCPU \i};
    \foreach \s in {0,1,2,3,4,5,6,7}{
      \pgfmathsetmacro{\x}{\s*0.78}
      \draw[run] (\x,\y-0.16) rectangle ++(0.72,0.32);}}
  \node[note, text=greenhl!55!black] at (6.5,2.78)
       {both vCPUs almost always\\find a free core};

  % ---- 8-vCPU guest: waits ----
  \node[font=\scriptsize\bfseries, text=primary, anchor=west] at (0,1.9)
       {The same work given 8 vCPUs};
  \foreach \i in {0,...,7}{
    \pgfmathsetmacro{\y}{1.4-\i*0.3}
    \node[lab] at (-0.15,\y) {vCPU \i};
    \foreach \s in {0,...,7}{
      \pgfmathsetmacro{\x}{\s*0.78}
      \pgfmathsetmacro{\r}{mod(\i+\s,3)}
      \ifdim\r pt>0.5pt
        \draw[run] (\x,\y-0.11) rectangle ++(0.72,0.22);
      \else
        \draw[wait] (\x,\y-0.11) rectangle ++(0.72,0.22);
      \fi}}
  \node[note, text=accent!70!black] at (6.5,0.1)
       {red is \emph{ready}: the vCPU had work\\
        and no core was free. More vCPUs\\
        means more chances to wait, and the\\
        guest's spinlocks make waiting\\
        cost its siblings too.};

  \draw[-{Stealth[length=2mm]}, neutral!60, line width=0.9pt]
        (0,-1.25) -- (6.3,-1.25);
  \node[font=\scriptsize\itshape, text=neutral, anchor=west] at (0,-1.6)
       {time};
  \node[rounded corners=2pt, draw=greenhl, fill=greenhl!15, inner sep=2pt,
        font=\tiny, text=greenhl!55!black] at (7.5,-1.25) {running};
  \node[rounded corners=2pt, draw=accent, fill=accent!12, inner sep=2pt,
        font=\tiny, text=accent!70!black] at (9.0,-1.25) {ready (waiting)};
\end{tikzpicture}}%
{Why a bigger guest can be a slower guest. Right-sizing downward is the most
common successful intervention in a contended cluster, and the one
administrators are most reluctant to make.}{vcpu-scheduling}

\begin{alertbox}[The rule that follows: size for what is used, not what is asked for]
Users ask for vCPUs the way they ask for disk --- more is safer. It is not. A
guest given eight vCPUs and using two has not merely wasted six; it has made
itself harder to schedule, raised its own ready time and raised everyone
else's.

\smallskip
The operational rule is to allocate what measurement justifies and add more
when ready time and utilisation both say so. \chref{performance} gives the
procedure; the point here is that the arithmetic of overcommit does not capture
it, because two guests with the same total vCPU count behave differently
according to how those vCPUs are grouped.
\end{alertbox}

\section{Memory, and Why It Is the Constraint}

\chref{what} asserted that memory, not processor time, usually decides how many
guests fit. The reason is that the two resources are elastic in different ways.

A processor can be oversubscribed because \emph{time} is divisible: two guests
each wanting a core can have it alternately, and the result is slower, not
wrong. Memory cannot be divided in time. A page either holds a guest's data or
it does not, and a guest denied a page it is using does not slow down slightly
--- it swaps, and swapping is a hundredfold penalty, not a percentage.

So hypervisors oversubscribe memory cautiously, using four techniques of
increasing desperation.

\dsfig{%
\begin{tikzpicture}[font=\scriptsize,
  rung/.style={rounded corners=3pt, line width=1pt, align=left,
               text width=5.5cm, inner sep=5pt, font=\scriptsize},
  cost/.style={rounded corners=2pt, line width=0.8pt, align=center,
               minimum width=2.1cm, minimum height=0.62cm, font=\scriptsize},
  hd/.style={font=\scriptsize\bfseries}
]
  \node[rung, draw=greenhl, fill=greenhl!7, text=greenhl!45!black] (s1) at (0,3.3)
       {\textbf{1. Transparent page sharing}\\
        Identical pages across guests are stored once, copy-on-write. Forty
        guests running the same operating system share most of their kernel.};
  \node[rung, draw=secondary, fill=secondary!7, text=secondary!75!black] (s2) at (0,1.5)
       {\textbf{2. Ballooning}\\
        A driver inside the guest allocates memory and tells the hypervisor it
        may take those pages. The \emph{guest} chooses what to give up, using
        its own page replacement.};
  \node[rung, draw=goldhl, fill=goldhl!7, text=goldhl!45!black] (s3) at (0,-0.3)
       {\textbf{3. Memory compression}\\
        Pages that would be swapped are compressed into a cache in RAM
        instead. Slower than memory, far faster than disk.};
  \node[rung, draw=accent, fill=accent!7, text=accent!65!black] (s4) at (0,-1.95)
       {\textbf{4. Hypervisor swapping}\\
        Guest pages are written to disk by the hypervisor, which does not know
        which ones matter. The last resort, and a cliff.};

  \node[cost, draw=greenhl, fill=white, text=greenhl!45!black] at (4.9,3.3)
       {free};
  \node[cost, draw=secondary, fill=white, text=secondary!75!black] at (4.9,1.5)
       {cheap};
  \node[cost, draw=goldhl, fill=white, text=goldhl!45!black] at (4.9,-0.3)
       {noticeable};
  \node[cost, draw=accent, fill=accent!10, text=accent!65!black] at (4.9,-1.95)
       {\textbf{catastrophic}};

  \draw[-{Stealth[length=2.4mm]}, neutral!60, line width=1.1pt]
        (6.4,3.8) -- (6.4,-2.5);
  \node[font=\scriptsize\itshape, text=neutral, anchor=west, align=left]
       at (6.6,0.6) {increasing\\pressure,\\increasing\\damage};

  \node[font=\scriptsize\itshape, text=accent!65!black, anchor=north west,
        align=left] at (-2.9,-3.0)
       {The gap between steps 3 and 4 is the one that matters: the first three
        degrade gracefully and the fourth does not.\\
        A host reaching step 4 has been oversubscribed, not tuned.};
\end{tikzpicture}}%
{The reclamation ladder. A hypervisor works down it as pressure rises; an
administrator's job is to ensure it never reaches the bottom
rung.}{reclamation-ladder}

\begin{conceptbox}[Why Ballooning Works, and Why It Feels Wrong]
Ballooning looks like a trick, and students reasonably ask why the hypervisor
does not simply take the pages.

\smallskip
Because it does not know which ones to take. From outside, every guest page
looks alike; the hypervisor cannot tell a hot database buffer from a page of
cache the guest would discard without noticing. Taking the wrong page costs a
disk read in the guest, and the hypervisor would be guessing.

\smallskip
The guest kernel, by contrast, has spent decades getting good at exactly this
question. So the balloon driver asks the guest for memory in the one language
the guest understands --- it allocates it, as any program would --- and the
guest runs its own page replacement to decide what to give up. The hypervisor
then takes the pages backing the balloon, which it knows are not in use.

\smallskip
The consequence worth remembering: \textbf{ballooning requires the guest to
cooperate}. A guest without the balloon driver installed cannot be reclaimed
from gracefully, and pressure on that host goes straight from step~1 to
step~3. Installing guest tools is not cosmetic.
\end{conceptbox}

\begin{worked}[Can This Host Take Four More Guests?]
A host of 2 $\times$ 32 cores and 512 GB currently runs 28 guests, each
configured with 4 vCPUs and 16 GB. Measured over a month: mean \%RDY 4\%,
95th percentile 9\%; mean active memory 6.2 GB per guest. A team asks for four
more identical guests. Decide.

\smallskip
\textbf{Step 1 --- CPU overcommit now and after.}
\[ \sigma_{\text{cpu}} = \frac{28 \times 4}{64} = \frac{112}{64}
   = \mathbf{1.75} \qquad\text{after: } \frac{32 \times 4}{64}
   = \frac{128}{64} = \mathbf{2.00} \]
Ready time is already at 9\% at the 95th percentile, which is the edge of the
acceptable band. Raising oversubscription by 14\% will not improve it.

\smallskip
\textbf{Step 2 --- memory \emph{allocated}.}
\[ 28 \times 16 = 448 \text{ GB} \qquad\text{after: } 32 \times 16
   = \mathbf{512 \text{ GB}} \]
The host has 512 GB \emph{in total}, and the hypervisor itself needs some ---
call it 6\%, about 32 GB, for its own code, the per-guest overhead and the page
tables of \chref{hardware}. So the request allocates 512 GB of 480 GB usable:
\textbf{32 GB short} before a single guest does anything.

\smallskip
\textbf{Step 3 --- memory \emph{active}.} Today the guests actually touch
\[ 28 \times 6.2 = 173.6 \text{ GB} \]
of 448 GB allocated --- 39\%. After the change, $32 \times 6.2 = 198.4$ GB,
still far below 480 GB. So on \emph{today's behaviour} the memory fits
comfortably.

\smallskip
\textbf{Step 4 --- which of Steps 2 and 3 should decide?} This is the real
question, and the answer is: Step 3 for the normal case, Step 2 for the bad
one. Overcommitting memory is safe exactly as long as the guests do not all
become active together --- and the events that make them active together
(a patch window, a backup run, a traffic spike, a failover from another host)
are correlated by construction.

\smallskip
If all 32 guests touched their full allocation, demand would be 512 GB against
480 GB usable: 32 GB short, about 1 GB per guest to be reclaimed. Page sharing
and ballooning would absorb that --- rungs 1 and 2 --- so the host would not
swap. That is survivable.

\smallskip
\textbf{Step 5 --- the recommendation.} \textbf{Refuse the request as stated,
and make a counter-offer.} The binding problem is CPU, not memory: 9\% ready at
the 95th percentile is the symptom of a host already at its limit, and
$\sigma_{\text{cpu}} = 2.0$ will make it worse for all 32 guests, not just the
new four.

\smallskip
The counter-offer is to right-size. If the existing 28 guests are, like most,
using two vCPUs of their four, reducing them to 2 vCPUs gives
$\sigma_{\text{cpu}} = (28 \times 2 + 4 \times 4)/64 = 72/64 = 1.13$ --- and
the four new guests fit with ready time \emph{lower} than it is today.
\chref{performance} is how you establish whether that premise holds.
\end{worked}

\begin{alertbox}[Do not read a memory overcommit ratio without an active figure]
$\sigma_{\text{mem}} = 1.25$ tells you almost nothing on its own. If the guests
are using a third of what they were given, it is comfortable. If they are using
all of it, the host is already swapping.

\smallskip
Allocation is a promise; active memory is a fact. Plan against the fact, size
the safety margin against the promise, and know which of the two you are
quoting whenever you quote a number.
\end{alertbox}

\section{NUMA}

A two-socket server does not have one pool of memory. Each socket has its own
memory controller and its own attached DIMMs; a core reaches its own socket's
memory quickly and the other socket's memory across an interconnect, more
slowly. This is \term{NUMA} --- non-uniform memory access --- and a remote
access typically costs something in the region of 1.5 to 2 times a local one in
latency, with lower bandwidth.

\dsfig{%
\begin{tikzpicture}[font=\scriptsize,
  sock/.style={rounded corners=3pt, line width=1.1pt, minimum width=5.0cm,
               minimum height=2.3cm},
  core/.style={rounded corners=1pt, line width=0.5pt, minimum width=0.52cm,
               minimum height=0.42cm, font=\tiny},
  mem/.style={rounded corners=2pt, line width=0.9pt, minimum width=4.4cm,
              minimum height=0.55cm, font=\scriptsize, align=center},
  note/.style={font=\scriptsize\itshape, align=center}
]
  \node[sock, draw=primary, fill=primary!5] at (0,0) {};
  \node[font=\scriptsize\bfseries, text=primary] at (0,0.95) {Socket 0};
  \foreach \i in {0,...,7}{
    \pgfmathsetmacro{\x}{-1.85+mod(\i,4)*0.62}
    \pgfmathsetmacro{\y}{0.45-floor(\i/4)*0.5}
    \node[core, draw=primary!60, fill=white, text=primary] at (\x,\y) {};}
  \node[mem, draw=primary, fill=primary!12, text=primary] at (0,-0.78)
       {256 GB local};

  \node[sock, draw=purplehl, fill=purplehl!5] at (6.4,0) {};
  \node[font=\scriptsize\bfseries, text=purplehl!70!black] at (6.4,0.95)
       {Socket 1};
  \foreach \i in {0,...,7}{
    \pgfmathsetmacro{\x}{4.55+mod(\i,4)*0.62}
    \pgfmathsetmacro{\y}{0.45-floor(\i/4)*0.5}
    \node[core, draw=purplehl!60, fill=white, text=purplehl] at (\x,\y) {};}
  \node[mem, draw=purplehl, fill=purplehl!12, text=purplehl!70!black]
       at (6.4,-0.78) {256 GB local};

  \draw[{Stealth[length=2mm]}-{Stealth[length=2mm]}, accent, line width=1.1pt]
        (2.55,0.1) -- (3.85,0.1);
  \node[note, text=accent!70!black] at (3.2,0.62) {interconnect};
  \node[note, text=accent!70!black] at (3.2,-0.4) {1.5--2$\times$\\the latency};

  \node[rounded corners=2pt, draw=greenhl, fill=greenhl!10,
        text=greenhl!45!black, font=\scriptsize, align=left, inner sep=4pt,
        text width=5.3cm] at (-1.0,-2.5)
       {\textbf{A guest that fits one socket}\\
        16 vCPU, 200 GB. Placed entirely on socket 0: every access local.};
  \node[rounded corners=2pt, draw=accent, fill=accent!8,
        text=accent!65!black, font=\scriptsize, align=left, inner sep=4pt,
        text width=5.3cm] at (5.4,-2.5)
       {\textbf{A guest that does not}\\
        40 vCPU, 400 GB. Spans both. Roughly half its accesses cross the
        interconnect --- unless the guest is told, and can act on it.};
\end{tikzpicture}}%
{NUMA on a two-socket host. The green guest is a sizing decision, not a tuning
option: once a guest is wider than a socket, no amount of hypervisor
configuration makes its memory local again.}{numa}

\begin{conceptbox}[vNUMA: Telling the Guest the Truth]
A guest that must span sockets can still be scheduled well, but only if it
knows. \term{vNUMA} exposes the host's NUMA topology to the guest, so the guest
kernel's own NUMA-aware scheduler and allocator can keep each thread's memory
on the node where the thread runs. Databases and JVMs do this competently.

\smallskip
Two consequences follow, and both are commonly missed. First, vNUMA is usually
enabled automatically only above a threshold --- historically eight or nine
vCPUs --- so a guest just under it that still spans sockets gets no topology at
all. Second, the topology is fixed when the guest boots: migrating it
(\chref{migration}) to a host with a different socket layout leaves the guest
optimising for a machine it is no longer on.

\smallskip
The rule that survives all of this: \textbf{size guests to fit inside a NUMA
node wherever you can}. It is the cheapest performance decision in this chapter
and it is made at provisioning time, where nobody is looking.
\end{conceptbox}

\begin{pitfall}
\begin{tight}
  \item \textbf{Adding vCPUs to a slow guest.} If the guest is slow because the
        host is contended, more vCPUs raise its ready time and everyone
        else's. Measure \%RDY before changing anything.
  \item \textbf{Quoting an overcommit ratio with no active-memory figure.}
        Allocation is a promise and activity is a fact; a ratio without the
        second is not a capacity statement.
  \item \textbf{Treating the four reclamation techniques as interchangeable.}
        Page sharing is free and ballooning is cheap; swapping is a cliff,
        not a slope. A host that swaps is misconfigured, not merely busy.
  \item \textbf{Running guests without the balloon driver and then
        overcommitting.} Without guest cooperation the hypervisor skips the
        graceful rung entirely.
  \item \textbf{Ignoring NUMA until there is a problem.} It is decided when the
        guest is sized. A 40-vCPU guest on a 32-core socket cannot be tuned
        back into locality.
  \item \textbf{Disabling page sharing because of side-channel advice, without
        noticing the memory cost.} Sharing between guests was restricted by
        default after 2014 for good reasons, but on a host of forty identical
        guests it was recovering a substantial fraction of memory. Know which
        trade you have made.
\end{tight}
\end{pitfall}

\begin{lab}[Making a Host Hurt, on Purpose]
Drives one host into CPU contention and then into memory pressure, and watches
each symptom appear. Use a laptop you are not relying on; it will become
unresponsive at step 5. Expect forty minutes.

\begin{lstlisting}[language=bash]
#!/usr/bin/env bash
# ch06_pressure.sh -- ready time and the reclamation ladder, observed.
set -euo pipefail
G () { ssh -o StrictHostKeyChecking=no "ubuntu@$1" "$2"; }

# --- 1. How many physical cores are we actually working with? ----------
nproc ; lscpu | grep -E 'Socket|Core|NUMA'

# --- 2. Build four guests, each with as many vCPUs as the host has. ----
# Deliberate oversubscription: 4 guests x nproc vCPUs on nproc cores.
for n in 1 2 3 4; do
  sudo qemu-img create -f qcow2 -b /var/lib/libvirt/images/vss1.qcow2 \
       -F qcow2 "/var/lib/libvirt/images/p$n.qcow2" 10G
  sudo virt-install --name "p$n" --memory 2048 --vcpus "$(nproc)" \
    --disk "/var/lib/libvirt/images/p$n.qcow2",format=qcow2 \
    --disk seed.iso,device=cdrom --os-variant ubuntu22.04 \
    --import --graphics none --noautoconsole
done
sleep 40

# --- 3. Baseline: one guest busy, three idle. --------------------------
ip_of () {
  sudo virsh domifaddr "$1" | awk '/ipv4/{split($4,a,"/"); print a[1]}'
}
IP1=$(ip_of p1)
G "$IP1" "sudo apt-get -qq install -y stress-ng sysstat"
G "$IP1" "stress-ng --cpu \$(nproc) --timeout 30s --metrics-brief"

# --- 4. Now make all four busy and read steal time in the guest. -------
# %st in the guest IS the ready time of section 6.1, seen from inside.
for n in 2 3 4; do
  IP=$(ip_of "p$n")
  G "$IP" "sudo apt-get -qq install -y stress-ng" || true
  G "$IP" "nohup stress-ng --cpu \$(nproc) --timeout 60s >/dev/null 2>&1 &"
done
sleep 5
G "$IP1" "mpstat 2 5 | tail -4"       # watch the %steal column climb

# --- 5. Memory pressure: ask for far more than the host has. -----------
# Watch the balloon move before anything swaps.
sudo virsh dommemstat p1
for n in 1 2 3 4; do
  G "$(ip_of "p$n")" "nohup stress-ng --vm 1 --vm-bytes 85% \
           --timeout 60s >/dev/null 2>&1 &" || true
done
sleep 20
sudo virsh dommemstat p1               # balloon / swap_in / swap_out
free -m ; vmstat 2 4                   # the host's own view

# --- 6. Right-size, and measure the same thing again. ------------------
# Two vCPUs each instead of nproc: the overcommit ratio drops fourfold.
for n in 1 2 3 4; do sudo virsh destroy "p$n"; \
  sudo virsh setvcpus "p$n" 2 --config --maximum --config; \
  sudo virsh setvcpus "p$n" 2 --config; sudo virsh start "p$n"; done
sleep 40
G "$IP1" "mpstat 2 5 | tail -4"        # %steal should be far lower

# --- 7. Tear down. -----------------------------------------------------
for n in 1 2 3 4; do
  sudo virsh destroy "p$n" || true
  sudo virsh undefine "p$n" --remove-all-storage || true
done
\end{lstlisting}

\textbf{What to notice.} Step 4's \texttt{\%steal} is the worked example's
\%RDY, measured from inside the guest, and step 6 shows that halving the vCPU
count made every guest \emph{faster}. That is the chapter's least intuitive
claim and you have now produced it yourself.
\end{lab}

\begin{checkpoint}
\begin{tightnum}
  \item A host of 48 cores runs 20 guests of 6 vCPUs each. Give the CPU
        overcommit ratio. Mean \%RDY is 14\%. State what you would do and what
        you would measure first.
  \item Explain why the hypervisor asks the guest to give up memory through a
        balloon driver rather than simply reclaiming pages itself.
  \item A guest is configured with 40 vCPUs and 400 GB on a host with two
        32-core sockets and 256 GB per socket. Name the problem, say why no
        hypervisor setting fixes it, and give the one change that would.
\end{tightnum}
\end{checkpoint}

\begin{chaptersummary}
\begin{tight}
  \item A vCPU is a thread scheduled onto a physical core. Ready time ---
        \%RDY, or steal time inside the guest --- measures how often it had
        work and no core. Under 5\% healthy, over 10\% oversubscribed.
  \item More vCPUs can make a guest slower: a guest's spinlocks assume its
        processors run together, and relaxed co-scheduling makes a wide guest
        wait. Right-sizing downward is the commonest effective fix.
  \item Processor time is divisible and memory is not, which is why memory
        usually decides guest density.
  \item The reclamation ladder: page sharing (free), ballooning (cheap,
        needs the guest driver), compression (noticeable), hypervisor swapping
        (a cliff). The job is to never reach the fourth.
  \item Overcommit ratios mean nothing without an active-memory figure.
        Allocation is a promise; activity is a fact.
  \item NUMA is decided at provisioning time. Size guests inside a node where
        you can; where you cannot, make sure vNUMA tells the guest the truth.
\end{tight}
\end{chaptersummary}

\begin{reviewq}
  \item \textbf{(MCQ)} Ready time measures: (a)~how long the guest took to
        boot; (b)~the fraction of time a vCPU was runnable but not scheduled;
        (c)~the hypervisor's own CPU use; (d)~time spent in VM exits.
  \item \textbf{(MCQ)} Ballooning requires: (a)~extended page tables; (b)~an
        IOMMU; (c)~a cooperating driver inside the guest; (d)~shared storage.
  \item \textbf{(Short)} Explain co-scheduling and why relaxed co-scheduling
        replaced the strict form.
  \item \textbf{(Short)} Give the four reclamation techniques in order and say
        what each costs the guest.
  \item \textbf{(Short)} Why is a memory overcommit ratio of 1.3 potentially
        safe while a CPU overcommit ratio of 1.3 is almost always safe?
  \item \textbf{(Applied)} A host of 2 $\times$ 24 cores and 384 GB runs 30
        guests of 4 vCPUs and 12 GB. Measured: \%RDY mean 7\%, p95 15\%; active
        memory 4.1 GB per guest. A team wants six more identical guests.
        Work through the same five steps as the chapter's worked example and
        give a recommendation with a counter-offer.
  \item \textbf{(Applied)} A database guest is reported as slow. It has 32
        vCPUs and 512 GB on a two-socket host with 24 cores and 384 GB per
        socket. List what you would measure, in order, and for each possible
        finding say what you would change.
\end{reviewq}
